\label{p3-20260826:chap:operador-angular-dictamen}
\index{operador angular}
\index{fase angular!salida espectral}
\index{cierre!pre-dimensional}

Las dos secciones anteriores han reunido tres piezas que antes aparecían
separadas: una identificación isométrica entre proyectores, una medida de
cociente y un cociclo positivo de acción. La presente sección estudia la consecuencia
espectral de esa ley y sitúa el resultado dentro de la arquitectura completa
de HMT\@.

La separación de estatutos es precisa. La identificación triple es un teorema
exacto relativo a la nueva regla variacional interna y a un ancla orientada.
La medida y el carácter positivo son teoremas exactos de \HTr--\HAdd; el defecto
es relativo además a \HDef, la capacidad a \HRes y el coste local a \HLoc. El valor
angular es una salida intervalarmente certificada del operador \Gtr. Su
proximidad a la coordenada angular conjugada canónica es extraordinaria, pero el intervalo de
la diferencia excluye la igualdad.

\section{Descomposición del operador angular en \(27\) sectores de \(72\) estados}

El lector de todos los caminos \Gtr posee \(1944\) estados microscópicos. La
acción de traslación gruesa de \((\mathbb Z/3\mathbb Z)^3\) permite la
descomposición
\[
 1944=27\cdot72.
\]
Cada bloque de \(72\) estados se parametriza por
\[
 (r,s,m,d)
 \in(\mathbb Z/3\mathbb Z)^2\times\{0,1\}\times D_4,
\]
donde \(m\) es la hoja y \(d\) la dirección anterior. Desde cada estado
parten cuatro continuaciones cardinales. El TRIT de la celda de llegada
decide la hoja siguiente.

Con la acción del \cref{p3-20260826:chap:medida-accion-positiva}, el peso de una arista
es
\[
 \exp[-g(e)-\lambda_*s(e)].
\]
La transformada de caracteres reduce el operador a matrices complejas
\(L_{k_a,k_b}\) de tamaño \(72\times72\). El programa enumera las
aristas de cada bloque; no recibe una matriz espectral precalculada. La
iteración de potencias se conserva como exploración, pero el resultado que
sigue ya no descansa en ella: los radios espectrales se aíslan mediante una
certificación racional completa.

Sean \(\rho_{01}\) y \(\rho_{12}\) los radios espectrales de los sectores
orientados \((0,1)\) y \((1,2)\). Definimos la \emph{fase de banda}
\[
 \boxed{y_{\rm band}=\log\rho_{12}-\log\rho_{01}.}
\]
Para
\[
 \lambda_*
 =1-\frac5{468}-\frac3{11}\Delta_4,
\]
la certificación da
\[
\begin{aligned}
1.460655120862011111
&<\rho_{01}<1.460655120862404024,\\
1.515812008270944328
&<\rho_{12}<1.515812008271195416,
\end{aligned}
\]
y, por propagación intervalar,
\[
 0.037066226434427529
 <y_{\rm band}<
 0.037066226434862172,
\]
\[
 \boxed{
 2.123738337168943^\circ
 <C_{\rm band}:=\frac{180}{\pi}y_{\rm band}
 <2.123738337193847^\circ.}
\]
La cantidad primaria es \(y_{\rm band}\), un logaritmo adimensional de cociente
espectral. Los grados aparecen después al aplicar la carta
\(180/\pi\). El operador produce primero una fase interna y sólo después
una representación angular arquimediana.

\subsection{Cálculo exacto del radio espectral del operador angular}

Los bloques son no normales. Por ello, un residuo pequeño
\(\|Lv-\mu v\|\) no basta para demostrar que \(\mu\) sea el autovalor de
mayor módulo. La prueba utilizada aquí consta de cuatro pasos.

Primero se observan doce filas estructuralmente nulas en cada bloque. Al
expandir el polinomio característico por ellas se obtiene un factor
\(z^{12}\) y un menor principal de orden sesenta. En éste aparecen otras
ocho filas estructuralmente nulas. Por tanto,
\[
 \det(zI_{72}-L)=z^{20}\det(zI_{52}-A),
\]
y todo autovalor no nulo queda contenido, con multiplicidad, en un bloque
reducido \(A\) de orden cincuenta y dos.

Segundo, el certificado proporciona matrices gaussianas racionales
\(S,R\), con denominador \(10^{16}\). El verificador calcula exactamente
\(E=I-RS\) y obtiene
\[
 \|E\|_\infty\le8.5664\times10^{-14}
 \quad\hbox{y}\quad
 \|E\|_\infty\le3.7628\times10^{-14}
\]
para los sectores \((0,1)\) y \((1,2)\), respectivamente. La serie de
Neumann demuestra entonces que \(S\) es invertible.

Tercero, si \(A_0\) es el punto medio racional del bloque intervalar, la
identidad
\[
 S^{-1}AS=(I-E)^{-1}RAS
\]
acota la distancia a \(C_0=RA_0S\) por
\[
 1.25125\times10^{-13}
 \quad\hbox{y}\quad
 5.7037\times10^{-14}.
\]
Los discos de Gershgorin de \(C_0\), ampliados por esas cotas, contienen
los discos de la similitud exacta. En cada bloque un disco queda separado
de los otros cincuenta y uno; el segundo teorema de Gershgorin le asigna
exactamente un autovalor. Su módulo inferior supera el módulo superior de
todos los demás discos. Ésa es la prueba de que el autovalor aislado es el
radio espectral.

Cuarto, \(\pi,e,\log\pi,e^{-1},e^{-\lambda_*}\) y \(\sqrt3\) no se
aceptan como decimales de biblioteca. El verificador los encierra mediante
la fórmula de Machin, las series factorial y logarítmica con resto, la
serie alternante de la exponencial y una horquilla racional cuadrática para
\(\sqrt3\). La verificación emplea únicamente enteros y fracciones; el
argumento queda cerrado por la reducción racional, la serie de Neumann y
la separación de Gershgorin expuestas en este apartado. Las expansiones
decimales reproducibles a profundidad diez mil se publican, como testigos
finitos independientes, en el
\hyperref[app:desarrollos-10000-rev9]{apéndice de expansiones}.

\section{Operador efectivo de banda \texorpdfstring{\(T_{\rm band}\)}{Tband} y aislamiento de la fase angular}

Sea
\[
 x=\frac{50\pi}{9}\alpha_{\rm HMT}
 =0.127362829012869966\ldots
\]
y sea
\[
 R=\begin{pmatrix}1&0\\0&-1\end{pmatrix}.
\]
El operador bicapa efectivo asociado exclusivamente a la fase de banda es
\[
 \boxed{T_{\rm band}=\exp(-xI+y_{\rm band}R).}
\]
En la base de sectores orientados,
\[
 T_{\rm band}
 =\begin{pmatrix}
 e^{-x+y_{\rm band}}&0\\
 0&e^{-x-y_{\rm band}}
 \end{pmatrix}.
\]
Se verifican
\[
 \det T_{\rm band}=e^{-2x},
 \qquad
 \frac{(T_{\rm band})_{11}}{(T_{\rm band})_{22}}
 =e^{2y_{\rm band}}.
\]
Por tanto,
\[
 -\frac9{100\pi}\log\det T_{\rm band}=\alpha_{\rm HMT},
 \qquad
 \frac{90}{\pi}
 \log\frac{(T_{\rm band})_{11}}{(T_{\rm band})_{22}}
 =C_{\rm band}.
\]

Las dos recuperaciones tienen estatutos distintos. La segunda devuelve la
salida producida por la brecha. La primera vuelve a leer el
\(\alpha_{\rm HMT}\) que entró en \(x\); demuestra reversibilidad de la
carta, no una segunda derivación independiente de \(\alpha\).

La forma exponencial separa dos funciones. El factor común \(e^{-x}\)
describe la contracción compartida de las dos hojas. Los factores
\(e^{\pm y_{\rm band}}\) distinguen orientaciones conservando el determinante. Como
\(I\) y \(R\) conmutan, el par
\[
 \left(\det T_{\rm band},
 \frac{(T_{\rm band})_{11}}{(T_{\rm band})_{22}}\right)
\]
reconstruye unívocamente \((x,y_{\rm band})\).

\section{Comparación posterior y escala de precisión}

La sección siguiente construye otra fase, la \emph{fase regional
regularizada} \(y_{\rm reg}\) ---denotada allí también por \(y_*\)---. Para
formular aquí una comparación con ese resultado posterior, anticipamos la
notación
\[
 C_{\rm reg}:=\frac{180}{\pi}y_{\rm reg}
 =2.123738338968646^\circ.
\]
Esta línea no constituye una segunda definición ni una entrada del operador
espectral: cita la conclusión del
\cref{p3-20260826:thm:contraangulo-regional}, cuya construcción se desarrolla en el
apartado siguiente.
No se identifican los dos pares de objetos:
\[
 (y_{\rm band},C_{\rm band})
 \ne
 (y_{\rm reg},C_{\rm reg}).
\]
La diferencia intervalar certificada de la ley lineal es
\[
 -1.799703\times10^{-9}\ {}^\circ
 <C_{\rm band}-C_{\rm reg}
 <-1.774799\times10^{-9}\ {}^\circ,
\]
con diferencia relativa encerrada por
\[
 -8.475\times10^{-10}
 <\frac{C_{\rm band}-C_{\rm reg}}{C_{\rm reg}}
 <-8.356\times10^{-10}.
\]
La ley que produce \(C_{\rm band}\) no recibe \(C_{\rm reg}\), CKM, Barbero,
masas ni datos SI\@. Esta independencia de objetivo hace de la proximidad un
candidato de reconocimiento fuerte y falsable. La certificación espectral
elimina la antigua laguna numérica; la diferencia intervalar, que excluye
cero, fija ahora el límite exacto: \(C_{\rm band}\) es una salida espectral
pre-dimensional distinta de la realización regional
\(C_{\rm reg}\).

El término adicional
\[
 -\frac{\Delta_4^2}{27}
\]
en \(\lambda\) produce
\[
 C_{\rm num}\approx2.12373833894104980{}^\circ,
\]
a \(2.76\times10^{-11}\) grados de \(C_{\rm reg}\). El entero \(27\) es interno,
pues cuenta los caracteres de \((\mathbb Z/3\mathbb Z)^3\), pero ese
hecho no fuerza el coeficiente del segundo momento. La aditividad exacta de
\HAdd excluye términos cuadráticos. La variante se conserva como ablación y no
como ley principal.

\section{Dependencia funcional del radio espectral respecto de los componentes del operador}

Todas las filas siguientes usan el mismo operador; sólo cambia una entrada
de la ley. Son diagnósticos numéricos de ablación; el intervalo certificado
precedente corresponde exclusivamente a la ley completa
\HTr--\HAdd--\HDef--\HRes--\HLoc (y, dentro de ella, el defecto positivo se fija en el
estrato \HTr--\HAdd--\HDef).
\[
\begin{array}{l|c|c}
\text{regla}&\lambda&C_{\rm num}\text{ en grados}\\ \hline
\text{autodual}&1&2.082755794784418\\
\text{sólo microfibra}&1-5/468&2.123621809570142\\
\text{traza }5/468+(3/11)\Delta_4
&0.989285913019141&2.123738337181414\\
4/468\text{ en lugar de }5/468
&0.991422665155894&2.115535228138553\\
3/12\text{ en lugar de }3/11
&0.989288440210566&2.123728626433667\\
5/243\text{ en lugar de }5/468
&0.979393542015659&2.161907060464779\\
5/729\text{ en lugar de }5/468
&0.993110963140488&2.109064222037112\\
7/729\text{, medida de semillas}
&0.990367478915522&2.119584299852700.
\end{array}
\]
Las ablaciones demuestran que la salida distingue la microfibra completa,
la dimensión reducida once y la medida de emisiones. No asignan por sí
solas una probabilidad de azar: tal probabilidad exigiría definir antes un
espacio de modelos alternativos y una medida sobre él.

\section{Composición tipada de los morfismos arquimediano, excepcional y dimensional}

La construcción completa puede escribirse como una sucesión de aplicaciones
con dominios visibles:
\[
\begin{aligned}
\Omega_{\rm seed}
&\xrightarrow{F}\mathcal U_6
\xrightarrow{J_F}H_U,\\
\mathscr R_\pi^{++}
&\xrightarrow{\beta_\pi}P_GH_G,\\
L_{\rm or}\otimes P_GH_G
&\xrightarrow{W_{GK}}P_3V_{11}
\xrightarrow{U_W}Q_3E_{30},\\
(\Pi_\pi^{(468)},\Delta_4P_3)
&\longmapsto D_{\rm switch}
\xrightarrow{\tau}\delta_*
\xrightarrow{\rm \HRes}\lambda_*,\\
\lambda_*
&\xrightarrow{\rm \HLoc}a_*
\xrightarrow{\mathcal L_{1,*}}y_{\rm band}
\xrightarrow{180/\pi}C_{\rm band},\\
(\alpha_{\rm HMT},y_{\rm band})
&\xrightarrow{\exp(-xI+y_{\rm band}R)}T_{\rm band}.
\end{aligned}
\]
Cada flecha declara lo que conserva:
\begin{itemize}
  \item \(F\) conserva la clase de emisión y deja en la fibra su
  multiplicidad de presentación;
  \item \(J_F\) compensa esa multiplicidad en la norma;
  \item \(\beta_\pi\) conserva la posición excepcional;
  \item \(W_{GK}\) conserva la representación orientada;
  \item \(U_W\) conserva incidencia y métrica;
  \item la traza olvida coordenadas y conserva rangos;
  \item el operador de transferencia suma caminos con el peso de la acción;
  \item la carta \(180/\pi\) representa en grados una fase interna.
\end{itemize}

\section{Diagrama de dependencias del operador angular y sus publicaciones}

\begin{longtable}{>{\raggedright\arraybackslash}p{.31\textwidth}>{\raggedright\arraybackslash}p{.22\textwidth}>{\raggedright\arraybackslash}p{.37\textwidth}}
\toprule
Objeto&Naturaleza matemática&Contenido demostrado\\
\midrule
\endfirsthead
\toprule
Objeto&Naturaleza matemática&Contenido demostrado\\
\midrule
\endhead
\(U_W=I/6:V_{11}\to E_{30}\)
&teorema exacto
&isometría, sobreyectividad, incidencia equivariante y unicidad positiva\\
\(P_3\leftrightarrow Q_3\)
&teorema exacto
&proyectores conjugados, rango tres y traza \(3/11\)\\
estrella \(R_{B_K}\)
&teorema exacto tipado
&espectro \(1^{[7]},(-1)^{[4]}\), índice \(3/11\) y obstrucción de conjugación\\
\(P_G\) en \(t=7\)
&teorema finito exacto
&filtración \(11\to3\to1\), proyector positivo y acción \(S_3\)\\
\(\mathscr R_\pi^{++}\leftrightarrow P_G\)
&teorema exacto
&bijección \(S_3\)-equivariante por posición excepcional\\
identificación isométrica orientada
&teorema relativo nuevo
&selector único por \(B_K,u_K\), Reynolds, polar y reducción \(240\to1\)\\
medida \(1/468\)
&teorema exacto de \HTr
&traza normalizada única y aplicación polar inducida por fibras\\
\(D_{\rm switch}\) y \(\delta_*\)
&teorema exacto relativo a \HTr--\HAdd--\HDef
&positividad, espectro y traza \(5/468+(3/11)\Delta_4\)\\
\(\lambda_*\)
&teorema interno relativo a \HTr--\HAdd--\HDef--\HRes
&capacidad residual obtenida del defecto mediante la regla \HRes\\
\(\mathcal A_*\)
&teorema interno relativo a \HTr--\HAdd--\HDef--\HRes--\HLoc
&cociclo aditivo de caminos con coste local \(g+\lambda_*s\)\\
\(T_{\rm band}\) y \(C_{\rm band}\)
&teorema computacional intervalar
&reducción \(72\to60\to52\), similitud racional, aislamiento de
Gershgorin y salida cuya distancia certificada está entre
\(1.774799\times10^{-9}\) y \(1.799703\times10^{-9}\) grados de
\(C_{\rm reg}\)\\
\(C_{\rm band}=C_{\rm reg}\)
&afirmación refutada para la ley lineal
&el intervalo certificado de la diferencia excluye cero\\
transición total
\(\mathcal X_{\mathrm{TPK}}^{\mathrm{enr}}\to \Gtr\)
&realizador global requerido
&\HTr--\HAdd--\HDef--\HRes--\HLoc define el lector relativo; la transición
total exige además una medida imagen dinámica sobre \(\Gtr\)\\
\bottomrule
\end{longtable}

\section{Acción predimensional en el centro de la doble proyección del estado genealógico}

La doble proyección parte del mismo estado enriquecido y produce dos
lecturas. En la rama arquimediana, los cilindros compatibles se refinan y
sus diámetros decrecen; la evaluación fija un valor límite y olvida parte
de la genealogía. En la rama excepcional, las incidencias se prolongan hacia
Witt, Golay, Niemeier--Leech, \(V^\natural\), el Monstruo y moonshine.

El nuevo resultado añade un objeto común más fuerte que una coincidencia de
cardinales. El módulo de once modos, su sector positivo de rango tres y la
traza de la acción pasan de dodecafase a Witt mediante \(U_W\). El sector orientado
\(t=7\) se incorpora mediante \(W_{GK}\). La microfibra de
\(\pi\) aporta el proyector regional \(\Pi_\pi^{(468)}\), y la medida uniforme del
catálogo fija su peso.

La acción positiva se sitúa así en el centro de las dos lecturas:
\[
 \text{región}
 \quad\oplus\quad
 \text{polarización orientada}
 \quad\longmapsto\quad
 \text{defecto positivo}.
\]
La evaluación arquimediana y la evaluación excepcional siguen siendo
aplicaciones diferentes; el entrelazador identifica un sector común, no
colapsa sus codominios.

\section{Compatibilidad del operador angular con los cilindros del continuo discreto}

Una historia HMT determina cilindros anidados
\[
 C_1\supset C_2\supset\cdots
\]
con prefijos compatibles. Cuando el lector arquimediano satisface
\[
 \operatorname{diam}(C_n)\longrightarrow0,
\]
la completitud de \(\mathbb R\) fija a lo sumo un valor en la intersección.
La nueva acción añade una estructura antes de ese límite: distingue el coste
de la ruta, la clase de emisión y la multiplicidad de presentaciones que el
valor final ya no muestra.

Éste es el sentido matemático preciso de afirmar que el continuo se lee como
proyección de genealogías discretas. El sistema de cilindros y sus
evaluaciones se prolonga a toda escala; el teorema global de cobertura
demuestra
\(\operatorname{ev}_{\mathbb R}(\Omega_{\infty,K})=K\) para cada compacto
\(K\subset\mathbb R\) y, al dirigir los intervalos, que la imagen es toda
\(\mathbb R\). El teorema finito de este apartado controla la acción angular
en cada nivel y se integra en esa cobertura global; no la reabre ni la
condiciona. El descenso algebraico de cada operación conserva su verificación
propia.

\section{Compatibilidad con el corredor excepcional, Moonshine y las pantallas de teoría M}

El entrelazador \(U_W\) aterriza en el módulo de incidencia de \(W_{12}\).
Desde el código ternario, las dos continuaciones excepcionales permanecen
separadas:
\[
 C_W\longrightarrow W_{12}
 \xrightarrow[\text{relativo a }(D,\ell)]{}W_{24},
\]
y
\[
 C_W\xrightarrow{\ \phi\ } N(A_2^{12})
\xrightarrow{\ v_\alpha\ }\Lambda_{24}
\longrightarrow V_\Lambda
\longrightarrow V^\natural
\longrightarrow\mathbb M
\longrightarrow\text{moonshine}.
\]
La acción positiva aporta a esta rama un sector transportado y una línea de
orientación. Las flechas desde Leech hasta moonshine son las construcciones
clásicas aplicadas al objeto finito obtenido. La acción del Monstruo se
realiza sobre \(V^\natural\); las palabras HMT no se identifican aquí con un
módulo directo para esa acción.

Las pantallas \(12\to11\to10\), \(26\to24\) y la involución T-dual forman
una interfaz matemática exacta con las dimensiones y dualidades usadas en
teoría M. La identificación física completa pertenece a Mecánica
Dimensional y se tipa en la segunda parte mediante el mapa adicional hacia
\(g_{MN}\), \(C_{MNP}\), \(\psi_M\), acción y supersimetría. No se reserva
para otra obra ni se deduce de la acción del Monstruo.

\section{Factorización espectral y aislamiento de los sectores dominantes}

La línea pre-dimensional queda organizada en cuatro estratos.
\begin{enumerate}
  \item \textbf{Núcleo finito.} APP, TRIT, el emisor TPK, el catálogo
  \(104\,976\to468\to243\), las cinco regiones de \(\pi\), la
  dodecafase, \(K\), \(\alpha\), Witt y la rama excepcional poseen
  construcciones explícitas y certificados.
  \item \textbf{Puente operatorial.} La incidencia construye \(U_W\);
  el sector orientado \(t=7\) construye \(P_G\); \(B_K\) y \(u_K\) seleccionan
  la identificación isométrica orientada; los tres proyectores de rango tres
  quedan identificados relativamente a esa regla.
  \item \textbf{Extensión de acción.} \HTr fija la medida \(1/468\), \HAdd
  fija el carácter aditivo y \HDef declara el acoplamiento
  \(\Delta_4P_3\). \HRes transforma el defecto en capacidad y \HLoc lleva la
  capacidad a las aristas. Las cinco cláusulas producen la ley local.
  \item \textbf{Consecuencia espectral.} \Gtr produce
  \(T_{\rm band}\) y una fase angular pre-dimensional certificada, con
  diferencia relativa entre \(8.356\times10^{-10}\) y
  \(8.475\times10^{-10}\) respecto de \(C_{\rm reg}\).
\end{enumerate}

Esto permite un cierre fuerte y tipado de esta ruta: el emisor finito de la
sección coinductiva infinita, la incidencia, la identificación de
proyectores, la medida y la acción quedan cerrados
respecto de las primitivas declaradas, de la regla variacional
\(\mathcal E_B\) y de las cláusulas de acción \HTr--\HAdd--\HDef--\HRes--\HLoc. \HAdd
selecciona la forma lineal; \HDef selecciona su segunda coordenada; \HRes y \HLoc
especifican cómo se lee como coste. El teorema global de
cobertura de todo \(\mathbb R\), la inducción desde una transición total de
\(\mathcal X_{\mathrm{TPK}}^{\mathrm{enr}}\) y la equivalencia física con
teoría M se formulan como extensiones de los resultados anteriores.

\section{Compatibilidad entre las realizaciones matricial, angular y reticular}

Las representaciones construidas satisfacen simultáneamente las condiciones
de compatibilidad siguientes:
\begin{enumerate}
  \item cinco regiones de \(\pi\), con tres positivas, una negativa y una
  transversal;
  \item \(P_G,P_3,Q_3\) como proyectores positivos y \(R_{B_K}\) como
  involución de orientación;
  \item el sector homogéneo orientado \(t=7\), no el cociente heterogéneo de \(t=8\);
  \item la línea \(L_{\rm or}\) en la identificación isométrica triple;
  \item el carácter relativo de la regla variacional que selecciona
  \(H_*,c_*,r_*\);
  \item \(5/468\) sobre emisiones y \(7/729\) sobre semillas, sin
  intercambiarlos;
  \item la linealidad del carácter como consecuencia de \HAdd y el coeficiente
  \(\Delta_4\) como contenido separado de \HDef;
  \item \((y_{\rm band},C_{\rm band})\) como salida espectral certificada,
  separado de \((y_{\rm reg},C_{\rm reg})\), la realización regional
  canónica de la sección siguiente;
  \item \(\alpha_{\rm HMT}\) como entrada de \(x\), de modo que la lectura
  del determinante es reversible y no independiente;
  \item las ramas \(W_{12}\to W_{24}\) y
  \(A_2^{12}\to\Lambda_{24}\) como construcciones separadas;
  \item el lector relativo \HTr--\HAdd--\HDef--\HRes--\HLoc distinguido de
  una transición total, cuya medida imagen dinámica constituye una
  construcción adicional.
\end{enumerate}

Estas condiciones no son notas editoriales. Forman parte de la estructura
lógica del resultado.
